% =====================================================================
%  CARNET D'ENTRAINEMENT — FICHE 6 : Liaisons et structures de Lewis
%  THÈME 5 — La résonance
%  Fiche TUTO (à lire avant les exercices)
%  Compilation : pdflatex
% =====================================================================
\documentclass[11pt,a4paper]{article}
\usepackage[utf8]{inputenc}
\usepackage[T1]{fontenc}
\usepackage{lmodern}
\usepackage[french]{babel}
\usepackage{amsmath,amssymb,mathtools}
\usepackage{icomma}
\usepackage[version=4]{mhchem}
\usepackage{siunitx}
\sisetup{output-decimal-marker={,}}
\usepackage{xcolor}
\usepackage{colortbl}
\usepackage{tikz}
\usepackage[most]{tcolorbox}
\usepackage{enumitem}
\usepackage{array}
\usepackage{booktabs}
\usepackage{multicol}
\usepackage{fancyhdr}
\usepackage[hidelinks]{hyperref}
\usetikzlibrary{arrows.meta,positioning,calc,shapes.geometric,shapes.misc,fit,backgrounds}
\usepackage[a4paper,margin=2cm,headheight=26pt,headsep=8pt]{geometry}

% -------- palette --------
\definecolor{primary}{RGB}{146,95,160}
\definecolor{primarydark}{RGB}{92,52,104}
\definecolor{defcol}{RGB}{0,114,178}
\definecolor{thmcol}{RGB}{0,140,120}
\definecolor{excol}{RGB}{0,135,90}
\definecolor{methcol}{RGB}{90,90,100}
\definecolor{attncol}{RGB}{200,40,55}
\definecolor{astucecol}{RGB}{198,93,9}
\definecolor{atomO}{RGB}{220,55,45}
\definecolor{atomC}{RGB}{60,60,70}
\definecolor{lonepair}{RGB}{120,94,200}
\definecolor{bondcol}{RGB}{60,60,70}
\definecolor{piecol}{RGB}{198,93,9}

% -------- boîtes --------
\tcbset{boxbase/.style={enhanced,breakable,boxrule=0.9pt,arc=2.5pt,
   left=3mm,right=3mm,top=2.2mm,bottom=2.2mm,
   fonttitle=\bfseries,coltitle=white,toptitle=0.6mm,bottomtitle=0.6mm}}
\newtcolorbox{defbox}[1][]{boxbase,colback=defcol!6,colframe=defcol,
   title={\textsc{À retenir}\ifx\relax#1\relax\relax\else\textmd{ : \itshape #1}\fi}}
\newtcolorbox{methbox}[1][]{boxbase,colback=methcol!8,colframe=methcol,sharp corners,
   title={\textsc{Méthode}\ifx\relax#1\relax\relax\else\textmd{ : \itshape #1}\fi}}
\newtcolorbox{exbox}[1][]{boxbase,colback=excol!6,colframe=excol,
   title={\textsc{Exemple}\ifx\relax#1\relax\relax\else\textmd{ : \itshape #1}\fi}}
\newtcolorbox{attnbox}[1][]{boxbase,colback=attncol!6,colframe=attncol,
   title={\textsc{Piège à éviter}\ifx\relax#1\relax\relax\else\textmd{ : \itshape #1}\fi}}
\newtcolorbox{astucebox}[1][]{boxbase,colback=astucecol!8,colframe=astucecol,
   title={\textsc{Astuce}\ifx\relax#1\relax\relax\else\textmd{ : \itshape #1}\fi}}

% -------- en-tête --------
\pagestyle{fancy}\fancyhf{}
\fancyhead[L]{\small\textcolor{primarydark}{\textbf{Fiche 6 · Thème 5 — La résonance}}}
\fancyhead[R]{\small\textcolor{primarydark}{\textsc{Tuto}}}
\fancyfoot[C]{\small\thepage}
\renewcommand{\headrulewidth}{0.8pt}
\renewcommand{\headrule}{\hbox to\headwidth{\color{primary}\leaders\hrule height \headrulewidth\hfill}}

\setlength{\parindent}{0pt}
\setlength{\parskip}{3pt}

\begin{document}

\begin{tcolorbox}[boxbase,colback=primary!12,colframe=primary,
   title={\Large Thème 5 — La résonance}]
\textbf{Objectif du thème.} Comprendre pourquoi une \emph{seule} structure de Lewis ne suffit
parfois pas, savoir écrire les \textbf{structures limites} reliées par la double flèche
$\longleftrightarrow$, décrire la molécule réelle comme un \textbf{hybride de résonance}, et surtout
ne \emph{jamais} confondre la résonance avec un équilibre chimique.
\end{tcolorbox}

\section*{1. Le problème : une formule unique qui ment}

Prenons l'ion carbonate \ce{CO3^2-}. L'expérience montre que ses \textbf{trois} liaisons \ce{C-O}
sont \emph{strictement identiques} (même longueur, intermédiaire entre une simple et une double
liaison) et que les trois oxygènes sont équivalents. Pourtant, si l'on suit les règles de Lewis, on
aboutit à une formule qui place \emph{une} double liaison \ce{C=O} sur \emph{un} oxygène précis et
deux liaisons simples \ce{C-O^-} sur les deux autres : cette formule unique introduit une différence
entre les oxygènes qui \textbf{n'existe pas} dans la réalité.

\begin{defbox}[résonance et hybride de résonance]
La \textbf{résonance} est la \textbf{délocalisation} d'électrons $\pi$ (des liaisons multiples) et de
doublets libres au sein d'un édifice, lorsqu'une seule structure de Lewis ne suffit pas à le décrire.
On écrit alors plusieurs \textbf{structures limites} (ou \emph{formes de résonance}), reliées par la
double flèche $\longleftrightarrow$. La molécule réelle n'est \emph{aucune} de ces formes : c'est un
\textbf{hybride de résonance}, une sorte de « moyenne » de toutes les structures limites.
\end{defbox}

\section*{2. Exemple guidé : l'ion carbonate \ce{CO3^2-}}

La double liaison \ce{C=O} peut aussi bien se placer sur l'un \emph{ou} l'autre des trois oxygènes.
Comme ces trois positions sont équivalentes, on obtient \textbf{trois} structures limites de même
énergie, qui contribuent à parts égales :

\begin{center}
\begin{tikzpicture}[scale=0.72,every node/.style={font=\small}]
  \tikzset{
    Oatom/.style={font=\bfseries},
    Catom/.style={font=\bfseries},
    sbond/.style={bondcol,line width=1.1pt},
    dbond/.style={bondcol,double,double distance=1.7pt,line width=0.6pt},
  }
  % ---- Form 1 : double en haut ----
  \begin{scope}[xshift=0cm]
    \coordinate (C) at (0,0);
    \coordinate (T) at (90:1.35);
    \coordinate (BL) at (210:1.35);
    \coordinate (BR) at (330:1.35);
    \draw[dbond] (C)--(T);
    \draw[sbond] (C)--(BL);
    \draw[sbond] (C)--(BR);
    \node[Catom] at (C){C};
    \node[Oatom] at (T){O};
    \node[Oatom] at (BL){O};
    \node[Oatom] at (BR){O};
    \node[piecol] at ($(BL)+(-0.42,-0.12)$){$\ominus$};
    \node[piecol] at ($(BR)+(0.42,-0.12)$){$\ominus$};
  \end{scope}
  \node at (2.55,0){$\longleftrightarrow$};
  % ---- Form 2 : double en bas-gauche ----
  \begin{scope}[xshift=5cm]
    \coordinate (C) at (0,0);
    \coordinate (T) at (90:1.35);
    \coordinate (BL) at (210:1.35);
    \coordinate (BR) at (330:1.35);
    \draw[dbond] (C)--(BL);
    \draw[sbond] (C)--(T);
    \draw[sbond] (C)--(BR);
    \node[Catom] at (C){C};
    \node[Oatom] at (T){O};
    \node[Oatom] at (BL){O};
    \node[Oatom] at (BR){O};
    \node[piecol] at ($(T)+(0.0,0.42)$){$\ominus$};
    \node[piecol] at ($(BR)+(0.42,-0.12)$){$\ominus$};
  \end{scope}
  \node at (7.55,0){$\longleftrightarrow$};
  % ---- Form 3 : double en bas-droite ----
  \begin{scope}[xshift=10cm]
    \coordinate (C) at (0,0);
    \coordinate (T) at (90:1.35);
    \coordinate (BL) at (210:1.35);
    \coordinate (BR) at (330:1.35);
    \draw[dbond] (C)--(BR);
    \draw[sbond] (C)--(T);
    \draw[sbond] (C)--(BL);
    \node[Catom] at (C){C};
    \node[Oatom] at (T){O};
    \node[Oatom] at (BL){O};
    \node[Oatom] at (BR){O};
    \node[piecol] at ($(T)+(0.0,0.42)$){$\ominus$};
    \node[piecol] at ($(BL)+(-0.42,-0.12)$){$\ominus$};
  \end{scope}
\end{tikzpicture}
\end{center}

\centerline{\footnotesize\itshape Les trois formes limites équivalentes de \ce{CO3^2-}
(charge totale $-2$ sur chaque forme : deux $\ominus$ portés par les oxygènes simplement liés).}

\begin{methbox}[écrire les formes limites d'un édifice]
\begin{enumerate}[leftmargin=1.7em,topsep=1pt,itemsep=1pt]
  \item Partir de la meilleure structure de Lewis (celle des charges formelles minimales).
  \item Repérer une double liaison \emph{voisine} d'un doublet libre (ou d'une charge) : c'est là que
        les électrons peuvent se \emph{délocaliser}.
  \item Faire « glisser » les électrons : la double liaison redevient simple d'un côté, un doublet
        libre devient une double liaison de l'autre. \textbf{Les atomes, eux, ne bougent pas.}
  \item Recommencer pour chaque position équivalente : le nombre de formes limites équivalentes
        égale le nombre de positions possibles pour la liaison multiple.
\end{enumerate}
\end{methbox}

\begin{exbox}[l'ozone \ce{O3}]
On écrit \ce{O=O-O <-> O-O=O} : deux formes limites. Dans la molécule réelle, les deux liaisons
\ce{O-O} sont \emph{identiques}, de longueur intermédiaire entre une simple et une double liaison —
c'est l'hybride des deux formes. La molécule reste \textbf{coudée} : la résonance ne change pas sa
géométrie.
\end{exbox}

\begin{attnbox}[$\longleftrightarrow$ n'est pas $\rightleftharpoons$]
La double flèche de résonance $\longleftrightarrow$ ne signifie \textbf{pas} que la molécule oscille
entre deux structures. Il n'y a \emph{pas} deux espèces qui s'interconvertissent (cela, c'est
l'équilibre chimique, noté $\rightleftharpoons$), mais \textbf{une seule} espèce, réelle et stable,
que notre langage de Lewis est incapable de dessiner d'un seul trait.
\end{attnbox}

\begin{astucebox}[les deux réflexes à garder]
\textbf{(1)} D'une forme limite à l'autre, \emph{seuls les électrons bougent} — jamais les atomes,
jamais la géométrie. \quad\textbf{(2)} $\longleftrightarrow$ relie les formes d'\emph{une} espèce
délocalisée ; $\rightleftharpoons$ relie \emph{deux} espèces distinctes d'un équilibre. Ne les
confonds jamais.
\end{astucebox}

\end{document}
